\section*{Chapter \ref{ch:g10:http-clients-servers} --- The Web: Clients, Servers and HTTP}

\begin{solution}{exo:g10:http-clients-servers:1}
\omterm{def:g10:http-clients-servers:url}{Scheme} \code{https}, \omterm{def:g10:http-clients-servers:url}{host} \code{maps.brenmoor.example}, no port in the address, path \code{/route}, \omterm{def:g10:http-clients-servers:url}{query string}
\code{from=3&to=17}, \omterm{def:g10:http-clients-servers:url}{fragment} \code{map}. With \code{https} and no port, the browser uses port $443$.
\end{solution}

\begin{solution}{exo:g10:http-clients-servers:2}
\code{201}: success (the resource was created). \code{301}: redirection (the resource has moved for good).
\code{403}: client error (the \omterm{def:g10:http-clients-servers:http}{request} is refused). \code{500} and \code{503}: server errors (the server failed, or
is not available for now). The \code{4xx} codes blame the \omterm{def:g10:http-clients-servers:http}{request}; the \code{5xx} codes blame the server.
\end{solution}

\begin{solution}{exo:g10:http-clients-servers:3}
\code{GET /search?title=the+lost+code HTTP/1.1}: the \omterm{def:g10:http-clients-servers:form}{form}'s \omterm{def:g10:http-clients-servers:form}{parameter} goes into the \omterm{def:g10:http-clients-servers:url}{query string}, and in a \omterm{def:g10:http-clients-servers:form}{form}'s
\omterm{def:g10:http-clients-servers:url}{query string} the spaces become \code{+}.
\end{solution}

\begin{solution}{exo:g10:http-clients-servers:4}
\textit{Visit number 1}, then \textit{Visit number 2}; after the \omterm{def:g10:http-clients-servers:cookie}{cookies} are deleted, the browser sends no \omterm{def:g10:http-clients-servers:cookie}{cookie}, so
the server counts from zero again: \textit{Visit number 1}. The jar holds \code{visits=1}.
\end{solution}

\begin{solution}{exo:g10:http-clients-servers:5}
Cut the address at \texttt{\#} and keep what comes before; keep what follows the first \code{?}; split it at each
\code{&}; split each pair at its first \code{=} (\code{exercises/query_params.py}). An address without \code{?} gives
an empty dictionary.
\end{solution}

\begin{solution}{exo:g10:http-clients-servers:6}
The empty line after the \omterm{def:g10:http-clients-servers:http-header}{headers}, which tells the client that the \omterm{def:g10:http-clients-servers:http-header}{headers} are over, and the \code{Content-Length}
\omterm{def:g10:http-clients-servers:http-header}{header}, which tells it how many bytes of body to read. Without them the client cannot know where the \omterm{def:g10:http-clients-servers:http-header}{headers} end or
where the body ends, so it keeps waiting for more (\code{exercises/build_response.py}).
\end{solution}

\begin{solution}{exo:g10:http-clients-servers:7}
The function would end without sending anything: the client would wait for a \omterm{def:g10:http-clients-servers:http}{response} that never comes, until it
gives up. \omterm{def:g10:http-clients-servers:http}{HTTP} pairs every \omterm{def:g10:http-clients-servers:http}{request} with one \omterm{def:g10:http-clients-servers:http}{response}, so even a wrong \omterm{def:g10:http-clients-servers:http}{request} must be answered, with an error code
that tells the client what went wrong.
\end{solution}

\begin{solution}{exo:g10:http-clients-servers:8}
The \omterm{def:g10:http-clients-servers:http}{request} line \code{GET /hello HTTP/1.1} has $19$ characters, $21$ bytes with its line end. \code{Host: 127.0.0.1}
takes $15 + 2 = 17$, \code{Connection: close} takes $17 + 2 = 19$, and the empty line $2$: $38$ bytes of \omterm{def:g10:http-clients-servers:http-header}{headers} and
empty line. That makes $59$ bytes in all, the value computed by the function \code{request_bytes} of
\code{exercises/build_response.py}.
\end{solution}

\begin{solution}{exo:g10:http-clients-servers:9}
The browser's history, the server's log of \omterm{def:g10:http-clients-servers:http}{request} lines, and every machine on the way, since \code{http} is not
encrypted (also anyone looking at the screen, and any page the browser reports the address to). Use
\code{method="post"}, so that the password travels in the body, and serve the site over \code{https}.
\end{solution}

\begin{solution}{exo:g10:http-clients-servers:10}
\texttt{caf\%C3\%A9\%20\%26\%20co}: \textit{é} is the \omterm{def:g10:text-encodings:utf8}{UTF-8} bytes \code{C3 A9}, the space \code{20} and \code{&}
\code{26}. In a \omterm{def:g10:http-clients-servers:url}{query string} \code{&} separates two \omterm{def:g10:http-clients-servers:form}{parameters}, so a raw \code{&} inside a value would cut it in two
and start a \omterm{def:g10:http-clients-servers:form}{parameter} named \textit{ co}.
\end{solution}

\begin{solution}{exo:g10:http-clients-servers:11}
In \code{do_GET}, for the path \code{/sum}, read \code{a} and \code{b} from \code{parse_qs(url.query)}, convert them
with \code{int}, and answer the sum; if a \omterm{def:g10:http-clients-servers:form}{parameter} is missing (\code{KeyError}) or not a whole number
(\code{ValueError}), answer \code{400} with a short message (\code{exercises/site_plus.py}). The tests ask for
\code{/sum?a=2&b=40} and get \code{42}, and for \code{/sum?a=2} and get \code{400}.
\end{solution}

\begin{solution}{exo:g10:http-clients-servers:12}
There are $5$ \omterm{def:g10:http-clients-servers:http}{requests}: the page, $1$ stylesheet and $3$ images. With a new connection each: $5 \times 2 = 10$ round
trips. With one connection reused: $1 + 5 = 6$. The browser learns the stylesheet's address only by reading the HTML,
which arrives at the end of the page's round trip.
\end{solution}

\begin{solution}{exo:g10:http-clients-servers:13}
On a \omterm{def:g10:http-clients-servers:http}{request} without a known \omterm{def:g10:http-clients-servers:cookie}{session} \omterm{def:g10:http-clients-servers:cookie}{cookie}, make an identifier with \code{secrets.token_hex(8)}, store an empty
list under it in a dictionary on the server, and send \code{Set-Cookie: session=...}; then add the item to that
visitor's list (\code{exercises/site_plus.py}). The visitor can see and change only a random identifier, which
cannot be guessed; the basket, and anything the server must trust, never leaves the server.
\end{solution}

\begin{solution}{exo:g10:http-clients-servers:14}
The client follows \code{/a}, \code{/b}, \code{/a}, \dots\ A careless client would loop for ever. \code{urllib}
remembers the addresses it was sent to: after the same address has come back a few times (nine \omterm{def:g10:http-clients-servers:http}{requests} in our
run), it stops with an error saying that the redirections would lead to an infinite loop. A client can count the
redirections of one \omterm{def:g10:http-clients-servers:http}{request} and give up past a limit, or stop when an address comes back.
\end{solution}

\begin{solution}{exo:g10:http-clients-servers:15}
A crawler reads the administrator's page, finds the links \code{/delete?book=1}, \code{/delete?book=2}, \dots, and
\omterm{def:g10:http-clients-servers:http}{requests} each one with \omterm{def:g10:http-clients-servers:methods}{GET}, since following a link is a \omterm{def:g10:http-clients-servers:methods}{GET}: every book is deleted. \omterm{def:g10:http-clients-servers:methods}{GET} must be safe. Deleting
should be a \omterm{def:g10:http-clients-servers:methods}{POST}, sent by a \omterm{def:g10:http-clients-servers:form}{form} button on a page that only a logged-in administrator can reach; a crawler does not
submit \omterm{def:g10:http-clients-servers:form}{forms}.
\end{solution}

\begin{solution}{pb:g10:http-clients-servers:1}
\textbf{1.} \code{style.css}, with a \code{link} tag, and \code{logo.svg}, \code{map.svg} and \code{clock.svg}, with
\code{img} tags.

\textbf{2.} $5$: the page and its $4$ resources.

\textbf{3.} The page: the other addresses are known only once its HTML has been read.

\textbf{4.} \code{200} for all five.

\textbf{5.} $425$, $92$, $204$, $203$ and $205$ bytes: $1129$ bytes of content.

\textbf{6.} $5 \times 130 = 650$ bytes of \omterm{def:g10:http-clients-servers:http-header}{headers}, so about $1779$ bytes; the \omterm{def:g10:http-clients-servers:http-header}{headers} are more than a third of it.

\textbf{7.} About $5 \times 400 = 2000$ bytes sent, more than the $1779$ received: for small resources, the \omterm{def:g10:http-clients-servers:http}{requests}
cost as much as the answers.

\textbf{8.} $1779 \times 8 = 14{,}232$ bits, at $10^7$ bits per second: about $1.4$~ms, against $50$~ms for one round
trip.

\textbf{9.} $10$ round trips, $500$~ms.

\textbf{10.} $6$ round trips, $300$~ms.

\textbf{11.} $2$ round trips for the connection and the page; then each of the two new connections takes $1$ round
trip to open and $2$ for its two resources, at the same time as the other: $2 + 3 = 5$ round trips, $250$~ms.

\textbf{12.} The round trips: the bytes take about $1.4$~ms, each round trip $50$~ms.

\textbf{13.} $6$ round trips in $1000$~ms: a latency of about $167$~ms, which is common on a distant or poor
mobile connection.

\textbf{14.} Still $1 + 5 = 6$ round trips: the $304$ answers carry no body, but each question still costs a round
trip.

\textbf{15.} $2$: the connection and the page.

\textbf{16.} The page changes (new books, opening hours, the results of a search): a copy kept for a day could show
stale content.

\textbf{17.} $6 - 2 = 4$ round trips per visit, so $40$ per day, $2$ seconds of waiting saved.

\textbf{18.} The library page takes $5$ \omterm{def:g10:http-clients-servers:http}{requests} and $1129$ bytes of content, about $1.8$~KB with the \omterm{def:g10:http-clients-servers:http-header}{headers}. At a
latency of $50$~ms it loads in $10$ round trips, $500$~ms, with a new connection per \omterm{def:g10:http-clients-servers:http}{request}, and in $6$ round trips,
$300$~ms, with one reused connection: the time goes into round trips, not bytes.
\end{solution}
